% Generated draft: requires mathematical and source review.
\documentclass[11pt]{article}
\usepackage[margin=25mm]{geometry}
\usepackage{fontspec}
\setmainfont{Latin Modern Roman}
\usepackage{amsmath,amssymb,mathtools,unicode-math}
\setmathfont{Latin Modern Math}
\usepackage{enumitem,needspace,xurl,hyperref}
\setlength{\parindent}{0pt}
\setlength{\parskip}{5pt}
\setlength{\emergencystretch}{3em}
\newcommand{\N}{\mathbb{N}}
\newcommand{\Z}{\mathbb{Z}}
\newcommand{\Q}{\mathbb{Q}}
\newcommand{\R}{\mathbb{R}}
\newcommand{\C}{\mathbb{C}}
\newcommand{\F}{\mathbb{F}}
\newcommand{\A}{\mathbb{A}}
\newcommand{\PP}{\mathbb{P}}
\newcommand{\E}{\mathbb{E}}
\newcommand{\eps}{\varepsilon}
\newcommand{\dd}{\,\mathrm{d}}
\begin{document}
\section*{UnsolvedMath ID: 800456}
\textbf{Source problem code:} catalogue entry
\subsection{Definitions and mathematical statement}
Let \(n\geq 1\) be an integer and let \(\F_2=\{0,1\}\) be the field with two elements. Write \(\mathbf{1}=(1,\ldots,1)\in\F_2^n\). A polynomial function of degree at most \(2\) means a function admitting an expression
\[
q(x)=a_0+\sum_{j=1}^n a_jx_j+\sum_{1\leq j<k\leq n}b_{jk}x_jx_k,
\qquad a_0,a_j,b_{jk}\in\F_2,
\]
with arithmetic in \(\F_2\). Its associated real-valued quadratic phase is \(x\mapsto(-1)^{q(x)}\).

The source calls its target the indicator of \(\mathbf{1}\), but specifies the range \(\{-1,1\}\). These descriptions require a convention that the record does not supply. Define the ordinary indicator by
\[
D_n(x)=\begin{cases}1,&x=\mathbf{1},\\0,&x\ne\mathbf{1},\end{cases}
\]
and its two signed encodings by
\[
A_{n,\sigma}(x)=\sigma\bigl(2D_n(x)-1\bigr),\qquad \sigma\in\{-1,1\}.
\]
The two signed encodings differ by a global sign, which can be absorbed into the real coefficients below. The ordinary indicator \(D_n\) is a distinct target; the source does not unambiguously reconcile its indicator terminology with its specified range.

\textbf{Problem, under the stated signed-range interpretation.} Fix either \(\sigma\in\{-1,1\}\). Suppose \(m\geq1\) is an integer, \(q_1,\ldots,q_m:\F_2^n\to\F_2\) have degree at most \(2\), and \(c_1,\ldots,c_m\in\R\) satisfy the exact pointwise identity
\[
A_{n,\sigma}(x)=\sum_{i=1}^m c_i(-1)^{q_i(x)}
\qquad\text{for every }x\in\F_2^n.
\]
Determine a lower bound on \(m\) as a function of \(n\), valid for every such representation. No distinctness condition on the \(q_i\), or restriction on the real coefficients, is supplied.

The background attributes to Hamed Hatami the conjecture that an exponential lower bound is necessary. In quantified form, this asks whether there exist constants \(a>0\) and an integer \(n_0\geq1\) such that, for every integer \(n\geq n_0\), every representation above satisfies
\[
m\geq 2^{an}.
\]

As supplied background, the source reports that Hatami can show \(m\geq n\), and states that if every \(q_i\) has degree at most \(1\), then \(m=2^n\) is necessary and sufficient, understood as a claim about the minimum number of terms. These are attributed source statements, not independently verified results here. The degree-one comparison requires a boundary qualification at \(n=1\) under the signed interpretation; the source supplies none. The record also describes the Constant-Degree Hypothesis of Barrington, Straubing, and Th\'{e}rien, attributed to 1990 and concerning finite fields, as a similar conjecture. Its formulation is not supplied, so no equivalence is asserted.
\subsection{Short English statement}
How many quadratic phases \((-1)^{q_i(x)}\), with arbitrary real coefficients, are needed to represent exactly the signed Boolean function that singles out \((1,\ldots,1)\)? The supplied source attributes to Hatami the conjecture that exponentially many terms are necessary. The indicator convention and a boundary qualification in the linear-phase comparison require review. The registry labels this problem open; that status is supplied, not independently verified.
\subsection{Sources}
\begin{itemize}
\item Local UnsolvedMath registry, permanent ID \(800456\), source problem code \texttt{catalogue entry}, title ``Analysis of Boolean Functions --- Uncertainty Principle for Quadratic Fourier Analysis.'' The supplied record has status \texttt{open}; its creation and update timestamps are both \texttt{2026-10-16T14:00:00Z}. These are registry data, not independent mathematical verification.
\item O'Donnell, \emph{Open Problems in Analysis of Boolean Functions} (2012), Problem 27, ``Uncertainty Principle for Quadratic Fourier Analysis.'' Supplied source URL: \url{https://arxiv.org/abs/1204.6447}. The record identifies the extraction method as \texttt{arxiv\_latex} and gives the access date \texttt{2026-10-16}. This is supplied provenance; the source was not independently visited or verified for this draft.
\item Attribution: the registry lists O'Donnell and 2012 in its proposer fields, whereas its background explicitly gives ``Proposed/Source: Hamed Hatami, 2011.'' The conjecture and reported lower bound are attributed to Hatami in that background; O'Donnell is the named source-list author.
\item Related background reference: Barrington, Straubing, and Th\'{e}rien, source label \texttt{[BST90]} (1990), concerning the Constant-Degree Hypothesis over finite fields. The supplied record gives no full bibliographic entry, URL, or precise statement for this reference.
\end{itemize}
\end{document}
